% =============================================================================
%                       WEEK 5: ETHICS AND MORALITY
% =============================================================================
\section{Week 5: Ethics and Morality}

% -----------------------------------------------------------------------------
%                              5.1 TOPIC SUMMARY
% -----------------------------------------------------------------------------
\subsection{Topic Summary}

\begin{keybox}[Learning Objectives]
By the end of this week, you will be able to:
\begin{itemize}
  \item Identify key \textbf{ethical challenges} posed by AI.
  \item Explain the main \textbf{philosophical theories of ethics}: deontology, consequentialism, virtue ethics.
  \item Describe \textbf{Moor's categorization} of moral agents.
  \item Compare \textbf{top-down} (deontic logic) and \textbf{bottom-up} (ML) approaches to implementing moral agents.
  \item Analyze \textbf{utilitarianism} and its critiques using trolley problem thought experiments.
  \item Discuss how ethical theories might be \textbf{combined} in AI implementations.
\end{itemize}
\end{keybox}

\separator

\subsubsection*{The Big Picture}

Ethics concerns what we \textit{ought} to do. As AI becomes more autonomous, we must ensure AI systems behave ethically. This week covers three major ethical theories---deontology (rules/duties), consequentialism (outcomes), and virtue ethics (character)---and how they inform the design of \textbf{artificial moral agents} (AMAs). Utilitarianism (a form of consequentialism) is explored in depth, including how human moral intuitions can diverge from rational utilitarian calculations.

\separator

\subsubsection*{What Examiners Like to Ask}

\begin{itemize}
  \item Compare and contrast deontology, consequentialism, and virtue ethics.
  \item What is Kant's Categorical Imperative? What are its critiques?
  \item Explain Moor's categorization of moral agents (ethical impact, implicit, explicit, full).
  \item How do top-down and bottom-up approaches to AMAs differ?
  \item What is utilitarianism? How is it defined?
  \item Analyze the trolley problem variations: What do they reveal about moral intuitions?
  \item Distinguish between ``means to an end'' and ``side effect'' in moral reasoning.
  \item What are the critiques of utilitarianism? How might they be addressed?
  \item How might ethical theories be combined in AI implementations?
\end{itemize}

\bigseparator

% -----------------------------------------------------------------------------
%                 5.2 OVERVIEW + CORE CONCEPTS + EXTENDED THEORY
% -----------------------------------------------------------------------------
\subsection{Overview + Core Concepts + Extended Theory}

\subsubsection*{Roadmap}
\begin{enumerate}
  \item Ethical challenges for AI
  \item Key theories of ethics: Deontology, Consequentialism, Virtue Ethics
  \item Moral relativism
  \item Moor's categorization of moral agents
  \item Implementing artificial moral agents (top-down vs.\ bottom-up)
  \item Utilitarianism in depth
  \item Trolleyology: Moral intuitions vs.\ utilitarian calculation
  \item Combining theories for AMAs
\end{enumerate}

\bigseparator

% --- Ethical Challenges for AI ---
\subsubsection{Ethical Challenges for Artificial Intelligence}

\begin{keybox}[Why Are We More Worried Now?]
\begin{enumerate}
  \item \textbf{Non-segregated environments:} AI (e.g., autonomous cars) now operates in open, unpredictable spaces, not just restricted areas with trained personnel.
  \item \textbf{Increased autonomy:} AI makes decisions independently, without human intervention.
  \item \textbf{Unpredictable contexts:} ML enables deployment in situations designers cannot fully anticipate.
  \item \textbf{Lack of transparency:} ML algorithms are ``black boxes''---we can't explain how they decide.
\end{enumerate}
\end{keybox}

\begin{exbox}[Examples of Ethical Challenges]
\begin{itemize}
  \item \textbf{Autonomous cars:} Brake failure $\to$ injure 3 workers or swerve and injure driver?
  \item \textbf{Autonomous weapons:} R\&D in US, UK, China, Russia, etc.
  \item \textbf{ML in decisions:} Insurance, recruitment, mortgages, criminal justice $\to$ bias, fairness, transparency.
  \item \textbf{Social robots:} Sex robots, elderly care robots $\to$ anthropomorphization, ethical treatment.
  \item \textbf{Social media / IoT:} Data privacy, surveillance, manipulation.
  \item \textbf{Filtering algorithms:} Polarizing beliefs, targeted advertising, restricting choices.
  \item \textbf{Medical AI:} Decision support, surgical robots, diagnosis.
  \item \textbf{Superintelligence:} Value loading problem (AI goals misaligned with human values).
\end{itemize}
\end{exbox}

\begin{defbox}[Ethics and AI: Key Questions]
\begin{enumerate}
  \item How should \textbf{we} use AI ethically? (User responsibility)
  \item How do we \textbf{design} AI so it cannot act unethically? (Hardcoded constraints)
  \item How do we design \textbf{autonomous} AI to \textit{choose} ethically in unforeseen situations?
  \item How do we remain \textbf{aware} of AI's ethical impact as it becomes embedded in society?
\end{enumerate}
\end{defbox}

\bigseparator

% --- Key Theories of Ethics ---
\subsubsection{Key Theories of Ethics and Morality}

\begin{center}
\renewcommand{\arraystretch}{1.4}
\begin{tabular}{@{} l p{9cm} @{}}
\toprule
\textbf{Theory} & \textbf{Core Idea} \\
\midrule
\textbf{Deontology} & Focus on \textbf{duties and rights}; morality = following rules/norms/laws. \\
\textbf{Consequentialism} & Focus on \textbf{outcomes}; choose actions with the best consequences. \\
\textbf{Virtue Ethics} & Focus on \textbf{character}; cultivate virtues (honesty, courage, etc.) $\to$ act ethically. \\
\bottomrule
\end{tabular}
\end{center}

\begin{exbox}[Helping Someone in Need]
\begin{itemize}
  \item \textbf{Consequentialist:} Helping maximizes well-being.
  \item \textbf{Deontologist:} Helping accords with ``Treat others as you would want to be treated.''
  \item \textbf{Virtue Ethicist:} Helping is what a charitable/benevolent person would do.
\end{itemize}
\end{exbox}

\separator

\paragraph*{Deontology}

\begin{defbox}[Deontology]
Concerned with choices that are morally \textbf{forbidden} (prohibited), \textbf{required} (obligated), or \textbf{permitted}. Guides what we \textit{ought} to do by following rules that encode duties and respect rights.

\textbf{Two approaches:}
\begin{itemize}
  \item Agent-focused: Duties of the agent to be moral.
  \item Rights-focused: Respect for rights of others.
\end{itemize}
\end{defbox}

\begin{exbox}[Deontological Rules]
\textbf{Specific rules:}
\begin{itemize}
  \item Religious: Ten Commandments, Hindu prohibition on beef.
  \item Secular: Laws against murder, theft; Asimov's Laws of Robotics.
\end{itemize}

\textbf{General principles:}
\begin{itemize}
  \item \textbf{The Golden Rule:} Treat others as you would like to be treated.
  \item \textbf{Kant's Categorical Imperative (CI):} Act only according to that rule by which you can at the same time will that it should become a \textit{universal law}.
\end{itemize}
\end{exbox}

\begin{defbox}[Kant's Categorical Imperative]
\textbf{Formulation 1:} Act only according to that rule by which you can \textit{will} that it should become a \textbf{universal law}.

\textbf{Example:} If you consider lying to get what you want, you must will that \textit{everyone} lies. But then no one believes anyone, and lying fails. So lying is not permitted---it requires making an exception for yourself.

\textbf{Formulation 2:} Act so that you treat humanity (yourself and others) always as an \textbf{end} and never simply as a \textbf{means}.

$\to$ Slavery is wrong because slaves are used as means to the slave owner's ends.
\end{defbox}

\begin{tipbox}[Critiques of Deontology]
\begin{enumerate}
  \item \textbf{Too general/vague:} Doesn't always guide action in specific situations.
  \item \textbf{Inflexible:} Absolute prohibitions; what about exceptions? (Can update rules, but then need to list all exceptions---computationally infeasible.)
  \item \textbf{No priority among rules:} When duties/rights conflict, no clear resolution. (E.g., abortion: right to life vs.\ right to choose.)
  \item \textbf{Harmful consequences:} Following rules may lead to bad outcomes. (E.g., ``A sadist is a masochist who follows the Golden Rule.'')
\end{enumerate}
\end{tipbox}

\separator

\paragraph*{Virtue Ethics}

\begin{defbox}[Virtue Ethics]
Morally good actions flow from the cultivation of \textbf{good character} (virtues): honesty, courage, compassion, wisdom, justice, etc.

\textbf{Origins:} Socrates, Aristotle.

\textbf{Learning:} What is good is learned through intuition, induction, experience---observing good people, being rewarded for good behavior.

$\to$ Resonates with \textbf{machine learning}: learn ethical behavior through experience and reinforcement, not explicit rules.
\end{defbox}

\begin{tipbox}[Critiques of Virtue Ethics]
\begin{itemize}
  \item Different cultures have different virtues; no objectively ``right'' list.
  \item Does not provide guidance on which actions are permitted/prohibited, only on what character to cultivate.
  \item In novel situations, virtues may not give definitive guidance (same problem for ML agents).
\end{itemize}
\end{tipbox}

\separator

\paragraph*{Consequentialism}

\begin{defbox}[Consequentialism]
Actions should be judged \textbf{only by their consequences}. Specify what is ``good,'' then choose actions that maximize the good.

\textbf{No rules or rights:} All that matters is outcomes. (E.g., lying to save an innocent life may be justified.)

\textbf{Fits AI/AIXI model:} Agent maximizes expected utility (= the good).
\end{defbox}

\bigseparator

% --- Moral Relativism ---
\subsubsection{Moral Relativism}

\begin{defbox}[Moral Relativism]
\textbf{Argument:}
\begin{enumerate}
  \item Morality = value judgments in a given society.
  \item Other societies have their own judgments; we have ours.
  \item \textbf{Conclusion:} We should not judge other cultures.
\end{enumerate}

\textbf{Critique:} The conclusion is itself a value judgment announced as a universal truth (``no matter what culture, you should not judge''). This is self-contradictory.
\end{defbox}

\bigseparator

% --- Moor's Categorization ---
\subsubsection{Moor's Categorization of Moral Agents}

\begin{center}
\renewcommand{\arraystretch}{1.3}
\begin{tabular}{@{} l p{9cm} @{}}
\toprule
\textbf{Category} & \textbf{Description} \\
\midrule
\textbf{Ethical Impact Agent} & Any machine whose operations increase or decrease good in the world. (E.g., robot camel jockeys replacing child jockeys; search ranking algorithms.) \\
\textbf{Implicit Ethical Agent} & Designed not to have negative ethical effects; behavior constrained by designers. No explicit representation of ethics. (E.g., pilot collision warning systems; robotic lawnmowers that can't harm.) \\
\textbf{Explicit Ethical Agent} & Can reason about the best action in ethical dilemmas, including \textbf{novel} situations. Uses explicit representations of ethics (e.g., deontic logic rules). \\
\textbf{Full Ethical Agent} & Has \textbf{moral agency}: understands ethical impacts, freely chooses, can justify moral judgments. $\to$ Morally responsible and accountable. \\
\bottomrule
\end{tabular}
\end{center}

\begin{tipbox}[Key Question]
Are these categories mutually exclusive? (E.g., an implicit ethical agent may also be an ethical impact agent.)
\end{tipbox}

\bigseparator

% --- Implementing AMAs ---
\subsubsection{Implementing Artificial Moral Agents (AMAs)}

\begin{center}
\renewcommand{\arraystretch}{1.3}
\begin{tabular}{@{} l p{4.5cm} p{5cm} @{}}
\toprule
\textbf{Approach} & \textbf{Description} & \textbf{Correspondence} \\
\midrule
\textbf{Top-Down} & Encode moral theory (rules) in agent; agent applies rules to act ethically. Logic-based, deontic. & Deontology \\
\textbf{Bottom-Up} & Agent learns ethical behavior through experience (reinforcement learning). Values implicit, not articulated. & Virtue Ethics \\
\textbf{Hybrid} & Combines top-down and bottom-up. & Multiple theories \\
\bottomrule
\end{tabular}
\end{center}

\separator

\paragraph*{Top-Down: Deontic Logic Implementations}

\begin{defbox}[Deontic Logic]
Uses modalities:
\begin{itemize}
  \item $Oq$ = $q$ is \textbf{obligatory}.
  \item $Pq$ = $q$ is \textbf{permitted}.
  \item $Fq$ = $q$ is \textbf{forbidden}.
\end{itemize}
\end{defbox}

\begin{exbox}[Autonomous Car Example (UK)]
\begin{itemize}
  \item $F(\text{drive\_on\_right})$ --- forbidden to drive on right.
  \item $O(\text{protect\_driver})$ --- obligated to protect driver.
\end{itemize}

\textbf{Scenario:} In a tunnel, collision imminent; only way to avoid is to swerve right.

\textbf{Conflict:} Must prioritize rules.

\textbf{Rule:} If collision imminent and risk of injury on left $>$ risk on right, then $O(\text{protect\_driver}) > F(\text{drive\_on\_right})$.

\textbf{Problem:} What if not in a tunnel? What about risk to others? What if there's a child in oncoming car? $\to$ Must encode \textit{all} exceptions---computationally infeasible.
\end{exbox}

\begin{thmbox}[The Core Problem]
For autonomous agents in complex, unpredictable environments:
\begin{itemize}
  \item Deontic rules must be \textbf{general} (can't anticipate all situations).
  \item Challenge: How to \textbf{interpret} general principles in specific situations?
  \item Related to \textbf{commonsense reasoning} problem (Week 4).
\end{itemize}
\end{thmbox}

\begin{exbox}[Asimov's Laws of Robotics]
\begin{enumerate}
  \setcounter{enumi}{-1}
  \item A robot may not injure humanity, or, by inaction, allow humanity to come to harm.
  \item A robot may not injure a human being or, through inaction, allow a human being to come to harm.
  \item A robot must obey orders given by human beings except where such orders conflict with the First Law.
  \item A robot must protect its own existence as long as this does not conflict with the First or Second Laws.
\end{enumerate}

\textbf{Problem:} What does ``harm'' mean? A literal-minded robot might interrupt a surgeon cutting a patient. Unless context is fully described, rules can lead to unintended consequences.
\end{exbox}

\separator

\paragraph*{Bottom-Up: Machine Learning Implementations}

\begin{itemize}
  \item Agent learns ethical behavior through \textbf{reinforcement learning}.
  \item Equates with \textbf{virtue ethics}: learning through experience, not explicit rules.
  \item Agent maximizes rewards by optimizing utility functions.
  \item \textbf{Problem:} In novel situations, no past experience to draw on.
\end{itemize}

\bigseparator

% --- Utilitarianism ---
\subsubsection{Utilitarianism}

\begin{defbox}[Utilitarianism]
\textbf{Utilitarianism = Consequentialism + Happiness as the measure of good.}

\[
\text{Utilitarianism} = \text{Impartial Maximization of Happiness}
\]

\begin{itemize}
  \item \textbf{Impartial:} Everyone's happiness counts equally (related to Golden Rule).
  \item \textbf{Maximization:} Choose actions that produce the greatest net happiness.
\end{itemize}

\textbf{Founders:} Jeremy Bentham, John Stuart Mill (18th--19th century).
\end{defbox}

\begin{keybox}[Why Happiness?]
\begin{enumerate}
  \item \textbf{Universal:} True for everyone $\to$ basis for universal morality.
  \item \textbf{Foundational:} All values we cherish (family, friendship, freedom, justice) are ultimately linked to happiness. Subtract their positive impact on experience, and their value is lost.
\end{enumerate}

\textbf{What is happiness?}
\begin{itemize}
  \item Subjective experience of well-being.
  \item Counterfactual test: If something we value were absent, would happiness decrease? Then it contributes to happiness.
\end{itemize}
\end{keybox}

\begin{thmbox}[Act vs.\ Rule Utilitarianism]
\begin{itemize}
  \item \textbf{Act Utilitarianism:} Apply utility calculation case-by-case; choose action that maximizes utility in that situation.
  \item \textbf{Rule Utilitarianism:} Actions are justified if they conform to \textit{rules} that, if included in our moral code, would maximize utility overall. (Combines deontology and utilitarianism.)
\end{itemize}
\end{thmbox}

\bigseparator

% --- Trolleyology ---
\subsubsection{Trolleyology: Moral Intuitions vs.\ Utilitarian Calculation}

\begin{defbox}[The Trolley Problems]
Thought experiments involving a trolley hurtling toward workers tied to a track. You know nothing about the workers (age, character, etc.). The question: What would you do?
\end{defbox}

\separator

\paragraph*{The Switch Dilemma (Standard)}

A trolley is heading toward 5 workers. At point A, you can steer the trolley onto a side track, killing 1 worker instead of 5.

\textbf{Response:} \textbf{87\%} would pull the switch (kill 1, save 5).

\textbf{Key features:}
\begin{itemize}
  \item No physical force (just pull a switch).
  \item Death of 1 is a \textbf{side effect}, not a means to an end. (Counterfactually: if the 1 worker weren't there, you'd still save the 5.)
\end{itemize}

\separator

\paragraph*{The Footbridge Dilemma}

You are on a footbridge. A large man with a backpack is next to you. The only way to stop the trolley is to push him off the bridge; he dies but stops the trolley, saving 5.

\textbf{Response:} Only \textbf{31\%} would push the man.

\textbf{Why the difference?} (Same utilitarian calculation: kill 1, save 5.)
\begin{enumerate}
  \item \textbf{Physical force:} Actively pushing someone $\to$ emotional resistance to violence (evolved to preserve social cooperation).
  \item \textbf{Means to an end:} The man is used as a ``trolley stopper.'' His death is \textit{causally necessary} to save the 5. (Counterfactually: if you didn't push him, you wouldn't save the 5.)
\end{enumerate}

\separator

\paragraph*{Footbridge + Switch Dilemma}

You are on a footbridge. You can run to the end and pull a switch to divert the trolley, saving 5. But you must collide with the man, knocking him off the bridge (he dies).

\textbf{Response:} \textbf{81\%} would run to pull the switch.

\textbf{Why?}
\begin{itemize}
  \item Physical force: Yes.
  \item Means to an end: \textbf{No}---the man's death is a \textbf{side effect}. (Counterfactually: if he weren't there, you'd still save the 5.)
\end{itemize}

\separator

\paragraph*{Switch Dilemma (Variant)}

Same as standard, but this time, if the 1 worker were not on the side track, the trolley would loop back and kill the 5. So the 1 worker is used as a trolley stopper (means to an end).

\textbf{Response:} \textbf{81\%} would steer right.

\textbf{Why?}
\begin{itemize}
  \item No physical force.
  \item Means to an end: Yes, but without physical force, the aversion is reduced.
\end{itemize}

\separator

\begin{thmbox}[Key Insight: Physical Force + Means to End]
It is the \textbf{combination} of:
\begin{enumerate}
  \item \textbf{Physical force} (violence), and
  \item \textbf{Means to an end} (death is causally necessary for the goal)
\end{enumerate}
that triggers strong emotional resistance and interferes with utilitarian calculation.

\textbf{Why?}
\begin{itemize}
  \item \textbf{Physical force:} We evolved to avoid violence (disrupts cooperation, threatens survival).
  \item \textbf{Means to an end:} When planning actions, we cognitively ``see'' the causal chain. Harms that are \textit{means} are more salient $\to$ trigger emotional alarm. \textit{Side effects} are less salient.
\end{itemize}
\end{thmbox}

\begin{exbox}[International Law: Means vs.\ Side Effect]
\begin{itemize}
  \item \textbf{Illegal:} Intentionally bombing civilians to lower enemy morale (civilians = means to an end).
  \item \textbf{Legal:} Bombing a munitions factory, knowing civilians will die as ``collateral damage'' (side effect).
\end{itemize}

\textbf{Critique:} Both are foreseeable deaths. Why is there a moral distinction? The distinction between means and side effect is arguably \textbf{morally irrelevant}---but our intuitions are shaped by this cognitive bias.
\end{exbox}

\begin{thmbox}[Trolleyology Conclusion]
\begin{itemize}
  \item People's moral intuitions diverge from pure utilitarian calculation.
  \item This is due to evolved emotional responses (aversion to violence, salience of means vs.\ side effects).
  \item But what people \textit{emotionally think} is right doesn't make it right.
  \item \textbf{AIs} will not have these evolutionarily acquired biases. They could apply utilitarian calculations without emotional interference.
\end{itemize}
\end{thmbox}

\bigseparator

% --- Critiques of Utilitarianism ---
\subsubsection{Critiques of Utilitarianism}

\begin{center}
\renewcommand{\arraystretch}{1.3}
\begin{tabular}{@{} p{4cm} p{5.5cm} p{4.5cm} @{}}
\toprule
\textbf{Critique} & \textbf{Example} & \textbf{Response} \\
\midrule
\textbf{Not demanding enough} (leads to unethical choices) & 5 people need organ transplants. Kill 1 healthy person to harvest organs and save 5? & Which society is happier: one where this is legal, or one where it is not? The latter $\to$ utilitarianism doesn't actually endorse this. \\
\textbf{Too demanding} (requires unreasonable sacrifices) & Instead of buying a \$500 suit, donate to save children in poverty. & \textbf{Pragmatic utilitarianism:} Being a ``perfect'' utilitarian is anti-utilitarian (incompatible with human nature). Be a role model others can emulate $\to$ greater good. \\
\bottomrule
\end{tabular}
\end{center}

\begin{tipbox}[Utilitarianism and the Abortion Debate]
\begin{itemize}
  \item Pro-choice: Appeals to \textbf{right to choose}.
  \item Pro-life: Appeals to \textbf{right to life}.
\end{itemize}

Both are deontic rights $\to$ no clear resolution when they conflict.

\textbf{Utilitarian approach:} Appeal to \textbf{evidence} about consequences. Does legalizing abortion increase or decrease overall well-being? (E.g., reduced poverty, fewer backstreet abortions, reduced crime?) This can be studied empirically.

When moral issues are \textbf{settled} (e.g., slavery), express them as rights. But for \textbf{controversial} dilemmas, utilitarianism may offer a path forward.
\end{tipbox}

\bigseparator

% --- Combining Theories ---
\subsubsection{Combining Ethical Theories for AMAs}

\begin{thmbox}[Hybrid Approach]
Perhaps AMAs should combine all three ethical theories:

\begin{enumerate}
  \item \textbf{Machine Learning (Virtue Ethics):} Reward agents for ethically good behaviors. Equip them to do the right thing in \textbf{familiar} situations with little controversy.
  
  \item \textbf{Utilitarianism (Consequentialism):} For \textbf{novel, controversial} dilemmas, use utilitarian calculations (under bounded resources). AI can:
  \begin{itemize}
    \item Collect and analyze data/evidence.
    \item Model causal reasoning (superior to humans).
    \item Humans provide input on how actions impact well-being (subjective experience).
  \end{itemize}
  
  \item \textbf{Rules (Deontology):} Once controversies are settled, encode as \textbf{rules} for long-term planning. (= Rule Utilitarianism.)
\end{enumerate}
\end{thmbox}

\begin{tipbox}[Utilitarian AI: Computational Challenges]
\begin{itemize}
  \item \textbf{Causal reasoning:} Computationally intractable to model all effects of actions.
  \item \textbf{Happiness estimation:} Must estimate effects on all conscious beings.
  \item \textbf{Bounded rationality:} Aim for ``good enough'' calculations under limited resources (like weather forecasting---imperfect but improving).
\end{itemize}
\end{tipbox}

\bigseparator

% --- Summary ---
\subsubsection{Summary: Key Takeaways}

\begin{keybox}[Week 5 Key Points]
\begin{enumerate}
  \item \textbf{Ethical challenges for AI:} Autonomy, unpredictability, lack of transparency.
  \item \textbf{Deontology:} Duties, rights, rules (Golden Rule, Kant's CI). Critiques: inflexible, can't handle exceptions, no priority among rules.
  \item \textbf{Virtue Ethics:} Cultivate good character. Critiques: no guidance on specific actions, culture-dependent.
  \item \textbf{Consequentialism:} Judge actions by outcomes. Utilitarianism = impartial maximization of happiness.
  \item \textbf{Moral relativism:} Self-contradictory.
  \item \textbf{Moor's categories:} Ethical impact, implicit, explicit, full ethical agents.
  \item \textbf{Top-down AMAs:} Deontic logic; problems with commonsense reasoning and exceptions.
  \item \textbf{Bottom-up AMAs:} ML/reinforcement learning; problems in novel situations.
  \item \textbf{Trolleyology:} Physical force + means-to-end $\to$ interferes with utilitarian calculation. Our intuitions are biased; AI could reason without these biases.
  \item \textbf{Hybrid AMAs:} Combine ML (familiar situations), utilitarianism (novel dilemmas), rules (settled issues).
\end{enumerate}
\end{keybox}

\begin{defbox}[Key Terms]
\begin{itemize}
  \item \textbf{Deontology:} Ethics based on duties/rules.
  \item \textbf{Categorical Imperative:} Act only on rules you can will as universal laws.
  \item \textbf{Virtue Ethics:} Ethics based on character.
  \item \textbf{Consequentialism:} Ethics based on outcomes.
  \item \textbf{Utilitarianism:} Impartial maximization of happiness.
  \item \textbf{Act vs.\ Rule Utilitarianism:} Case-by-case vs.\ rule-based utility.
  \item \textbf{Means to an end vs.\ Side effect:} Causally necessary vs.\ incidental harm.
  \item \textbf{Implicit/Explicit/Full Ethical Agent:} Moor's categories.
  \item \textbf{Top-down/Bottom-up:} Rule-based vs.\ learning-based AMAs.
\end{itemize}
\end{defbox}

